\documentclass[10pt,a4paper]{amsart}
\usepackage[margin=19mm]{geometry}
\usepackage{amsmath,amssymb,mathtools}
\usepackage[scr=rsfs,cal=euler]{mathalfa}
\usepackage{xfrac}
\usepackage[unicode,hidelinks,bookmarks=false]{hyperref}
\newcommand{\ie}{i.e.\ }
\newcommand{\cf}{cf.\ }
\newcommand{\resp}{respectively\ }
\newcommand{\Subsection}{\subsection}
\providecommand{\nobreakdash}{\nobreak\hbox{-}}
\setlength{\emergencystretch}{2em}
\multlinegap0pt
\usepackage{amssymb}
\usepackage{mathtools}
\usepackage{bm}
\usepackage{siunitx}
\usepackage{float}
\usepackage[shortlabels]{enumitem}
\newtheorem{lemma}{Lemma}
\title{Intermittency induced by long memory under stochastic regime switching}
\author{Mauricio Herrera-Mar\'in}
\date{Source: arXiv:2605.00729v1 (CC BY 4.0)}
\begin{document}
\maketitle

\noindent\emph{Source and licence.} Intermittency induced by long memory under stochastic regime switching. Version arXiv:2605.00729v1 (CC BY 4.0), distributed by arXiv under the Creative Commons Attribution 4.0 International licence, http://creativecommons.org/licenses/by/4.0/, as recorded in the arXiv licence metadata for this exact version and shown on the abstract page https://arxiv.org/abs/2605.00729v1. Full licence notice: \texttt{licenses/arxiv-2605.00729-CC-BY-4.0.txt}.\par\smallskip

\noindent\textit{Editorial scope.} Counted unit: the article's lemma app:lem:generator (source0.tex lines 1803--1850). The article's complete original proof of it is delivered with it (source0.tex lines 1852--1899, one \texttt{proof} environment of 48 lines). No mathematics is written, repaired or re-derived: the statement and the proof are the article's own lines, in the article's own notation, and no retained line is substituted or re-flowed.  No further source-local context is needed: every cross-reference of every retained line resolves inside the two delivered spans.  Nothing was omitted from any retained span, so the package carries no omission record.  No retained line contains a citation, so the wrapper carries no bibliography and no bibliography entry.  No retained line contains a figure environment, an image-inclusion command or drawing code, so no image is delivered and the package declares no asset dependency.  Editorial formatting only, in addition: an AMS \texttt{amsart} preamble that restates the article's own declarations and macros used by the retained lines, namely the environments \texttt{lemma} on separate counters, together with the standard packages those lines need.  The retained environments print in this document as Lemma 1 in order of appearance, whereas the article prints the same items with the numbering of its own full document.  The complete native source of the article as distributed in the arXiv e-print for this exact version is bundled unchanged as \texttt{sources/199587/source0.tex}.
	\begin{lemma}[Infinitesimal generator for memory-dependent Lyapunov functionals]
		\label{app:lem:generator}
		Let $V:\mathcal{X}\times\mathcal{Z}\to\mathbb{R}_+$ be the Lyapunov functional
		defined by
		\begin{equation}
			\label{app:eq:Lyap-memory}
			V(u,z)
			=
			\|u\|_{\mathcal{X}}^2
			+
			\int_0^\infty \phi_z(s)\,\|u(t-s)\|_{\mathcal{X}}^2\,ds,
		\end{equation}
		where $\phi_z:\mathbb{R}_+\to\mathbb{R}_+$ is locally integrable, completely
		monotone, and $\int_0^\infty \phi_z(s)\,ds<\infty$.
		Then $V$ belongs to the domain of the infinitesimal generator $\mathcal{L}$ of
		the Markov process $(u(t),Z(t))$, and its action decomposes as
		\begin{equation}
			\label{app:eq:generator-decomposition}
			\mathcal{L}V(u,z)
			=
			\mathcal{L}_{\mathrm{det}}V(u,z)
			+
			\mathcal{L}_{\mathrm{jump}}V(u,z),
		\end{equation}
		where
		\begin{align}
			\label{app:eq:generator-det}
			\mathcal{L}_{\mathrm{det}}V(u,z)
			&=
			2\langle u,\mathcal{A}_z u\rangle_{\mathcal{X}}
			+
			2\Big\langle
			u,
			\int_0^\infty K_z(s)\,\mathcal{B}_z u(t-s)\,ds
			\Big\rangle_{\mathcal{X}}
			\nonumber\\
			&\quad
			-
			\int_0^\infty \phi_z'(s)\,\|u(t-s)\|_{\mathcal{X}}^2\,ds,
		\end{align}
		and
		\begin{equation}
			\label{app:eq:generator-jump}
			\mathcal{L}_{\mathrm{jump}}V(u,z)
			=
			\sum_{z'\neq z} q_{zz'}\big(V(u,z')-V(u,z)\big).
		\end{equation}
	\end{lemma}

	\begin{proof}
		We decompose the argument into deterministic and jump contributions.
		
		\medskip
		\noindent\emph{Step 1: Deterministic Volterra contribution.}
		Fix $z\in\mathcal{Z}$ and consider the evolution under the frozen regime
		$Z(t)\equiv z$.
		By resolvent-family theory for Volterra equations with completely monotone kernels,
		mild solutions $u(t)$ are continuous in $\mathcal{X}$ and locally absolutely
		continuous in time.
		Differentiating $\|u(t)\|_{\mathcal{X}}^2$ along trajectories yields
		\[
		\frac{d}{dt}\|u(t)\|_{\mathcal{X}}^2
		=
		2\langle u(t),\mathcal{A}_z u(t)\rangle_{\mathcal{X}}
		+
		2\Big\langle
		u(t),
		\int_0^t K_z(t-s)\,\mathcal{B}_z u(s)\,ds
		\Big\rangle_{\mathcal{X}}.
		\]
		For the memory term, differentiation under the integral sign and complete
		monotonicity of $\phi_z$ yield $\phi_z'\le 0$ in the sense of distributions and
		\[
		\frac{d}{dt}
		\int_0^\infty \phi_z(s)\,\|u(t-s)\|_{\mathcal{X}}^2\,ds
		=
		-
		\int_0^\infty \phi_z'(s)\,\|u(t-s)\|_{\mathcal{X}}^2\,ds.
		\]
		Combining both contributions gives \eqref{app:eq:generator-det}.
		
		\medskip
		\noindent\emph{Step 2: Jump contribution.}
		Since $V$ depends on the regime only through $z$, the jump part of the generator is
		\[
		\mathcal{L}_{\mathrm{jump}}V(u,z)
		=
		\sum_{z'\neq z} q_{zz'}\big(V(u,z')-V(u,z)\big),
		\]
		which is \eqref{app:eq:generator-jump}.
		
		\medskip
		\noindent\emph{Step 3: Domain considerations.}
		The assumed integrability of $\phi_z$ and the boundedness of the resolvent family
		ensure all terms are finite and measurable, hence $V\in\mathrm{Dom}(\mathcal L)$ and
		\eqref{app:eq:generator-decomposition} holds.
	\end{proof}

\end{document}
